\documentclass[12pt]{article}

\usepackage[T1]{fontenc}
\usepackage[protrusion=true,expansion=false]{microtype}
\usepackage[margin=1.25in, headheight=14.5pt]{geometry}
\usepackage{fancyhdr}
\pagestyle{fancy}
\fancyhf{}
\fancyhead[R]{\thepage}
\fancyhead[L]{\small\itshape A Factorization Criterion for Route Invariants}
\renewcommand{\headrulewidth}{0pt}
\fancypagestyle{plain}{\fancyhf{}\renewcommand{\headrulewidth}{0pt}\fancyfoot[C]{\thepage}}
\usepackage{amsmath,amssymb,mathtools,amsthm}
\usepackage{aliascnt}
\usepackage[unicode,hypertexnames=false]{hyperref}
\hypersetup{
  pdftitle={A Factorization Criterion for Route Invariants with Fixed Endpoint Data},
  pdfauthor={Chast K. Wolfe},
  pdfsubject={endpoint data; route invariants; transport; holonomy;
    curvature; quotient descent},
  pdfkeywords={endpoint data, route invariants, transport, holonomy,
    gauge theory, category theory, curvature, general relativity},
  hidelinks
}

\newtheorem{theorem}{Theorem}[section]
\newaliascnt{proposition}{theorem}
\newtheorem{proposition}[proposition]{Proposition}
\aliascntresetthe{proposition}
\newaliascnt{lemma}{theorem}
\newtheorem{lemma}[lemma]{Lemma}
\aliascntresetthe{lemma}
\newaliascnt{corollary}{theorem}
\newtheorem{corollary}[corollary]{Corollary}
\aliascntresetthe{corollary}
\newaliascnt{definition}{theorem}
\newtheorem{definition}[definition]{Definition}
\aliascntresetthe{definition}
\newaliascnt{principle}{theorem}
\newtheorem{principle}[principle]{Principle}
\aliascntresetthe{principle}
\newaliascnt{remark}{theorem}
\newtheorem{remark}[remark]{Remark}
\aliascntresetthe{remark}

\usepackage[nameinlink,noabbrev]{cleveref}

\crefname{theorem}{Theorem}{Theorems}
\crefname{proposition}{Proposition}{Propositions}
\crefname{lemma}{Lemma}{Lemmas}
\crefname{definition}{Definition}{Definitions}
\crefname{principle}{Principle}{Principles}
\crefname{corollary}{Corollary}{Corollaries}
\crefname{section}{Section}{Sections}

\newcommand{\calP}{\mathcal{P}}
\newcommand{\calR}{\mathcal{R}}
\newcommand{\calC}{\mathcal{C}}
\newcommand{\id}{\operatorname{id}}
\newcommand{\Mor}{\operatorname{Mor}}
\newcommand{\Ob}{\operatorname{Ob}}
\newcommand{\Hom}{\operatorname{Hom}}
\newcommand{\Aut}{\operatorname{Aut}}
\newcommand{\Hol}{\operatorname{Hol}}
\newcommand{\Data}{\mathsf{D}}
\newcommand{\citep}[1]{\cite{#1}}

\begin{document}

\title{A Factorization Criterion for Route Invariants with Fixed Endpoint Data}
\author{Chast K.~Wolfe\\[4pt]
\small Closure Research Initiative, United States\\[2pt]
\small Corresponding author:~\href{mailto:chast.wolfe@closureresearchinitiative.org}{\texttt{chast.wolfe@closureresearchinitiative.org}}}
\date{Version 3\\June 13, 2026}
\maketitle

\begin{abstract}
\noindent
A factorization criterion is proved for comparison data with fixed
endpoints. A comparison framework consists of endpoints, routes,
endpoint data, and route invariants. The endpoint map sends each route
to the data determined by its source and target. The main result is that
a route invariant is determined by endpoint data if and only if it is
constant on every endpoint fiber, equivalently if and only if it factors
through the endpoint quotient. Equivalently, endpoint closure holds
precisely for the invariants that descend to the endpoint quotient.
Endpoint data specify source, target, and any datum assigned to an
ordered pair; they do not include the route itself. Thus any invariant
separating routes with the same source, target, and endpoint datum is not
endpoint-determined. With composition, a basic obstruction is the direct
route versus a composite route. When a return route is available, the
discrepancy can be completed to a loop. In smooth transport geometry,
infinitesimal transport around an infinitesimal loop is represented by
curvature. The proof uses elementary quotient descent for the endpoint
map. Subsequent sections record instances in quotient semantics,
categories, networks, gauge transport, Wilson loops, holonomy,
differential geometry, and general relativity. The result is a descent
criterion, not a derivation of those theories. The formal obstruction is
that representative-, morphism-, path-, transport-, and loop-dependent
information is not endpoint-determined unless it is constant on endpoint
fibers.
\end{abstract}

\paragraph{Keywords.}
endpoint data; route invariants; quotient descent; holonomy;
curvature; comparison data

\paragraph{Conventions.}
All collections below may be read as sets, or as classes in a fixed
universe. When quotient notation is used, it denotes the ordinary set- or
class-quotient by the displayed equivalence relation. No topological or
smooth structure is assumed unless explicitly stated. The letters
$p,q,r$ denote endpoints; $\gamma,\eta$ denote routes or paths; $s,t$
denote source and target maps; $C$ denotes endpoint-comparison data; $e$
denotes the endpoint map; and $\pi_E$ denotes the endpoint quotient map.

\section{Endpoint Data and Route Invariants}
\label{sec:intro}

The basic operation considered here is passage from routes to the data
determined by their endpoints. This operation is a map from a finer
description to coarser data, analogous to a quotient map or a projection
to boundary data. A quantity descends along such a map exactly when it is
constant on its fibers. In the present setting, those fibers consist of
routes with the same endpoint data.

Write an endpoint comparison as
\[
  C(p,q),
\]
where $p$ and $q$ are endpoints. The symbol may denote equality,
equivalence, distance, causal relation, boundary value, quotient class,
or any other datum assigned to the ordered pair. The datum is determined
by endpoints. It contains neither a route between the endpoints, nor a
composition law for comparisons, nor equality of route-dependent values
for two routes from $p$ to $q$.

The issue is factorization through the endpoint map, equivalently
descent of the chosen invariant along endpoint projection. Endpoint
data determine exactly the invariants that are constant on endpoint
fibers. If an invariant distinguishes routes with the same source and
target, then it is not determined by endpoint data. Representation of
such an invariant requires more than endpoints.

Let $\calR$ be a class of routes and let
$\pi_E:\calR\to\calR/{\sim_E}$ send a route to its endpoint-data
class. A route invariant $I$ is determined by endpoint data if and only
if it factors through $\pi_E$. Equivalently, $I$ must be constant on every
fiber of $\pi_E$. This is the ordinary descent condition for a quotient
\citep{Munkres2000}.
In the closure program, this is the quotient-descent form of the
route/transport obstruction used in the monograph \citep{WolfeCSM}; no
monograph result is assumed in the proof below.

Several standard settings instantiate the same factorization condition.
For quotients, descent is constancy on equivalence classes
\citep{Munkres2000}. In category theory, object data do not determine
morphisms \citep{EilenbergMacLane1945,MacLane1998}. In networks, path
products need not be determined by endpoints. In gauge theory, Wilson
loops are gauge-invariant loop data rather than endpoint data
\citep{Wilson1974}.
In differential geometry, connections, holonomy, and curvature encode
transport around loops
\citep{Ehresmann1950,AmbroseSinger1953,KobayashiNomizu1963,Nakahara2003}.
In general relativity, the corresponding hierarchy is metric comparison,
Levi-Civita transport, curvature, and contraction, with curvature
downstream of transport \citep{Einstein1916,Wald1984,Malament2012,Carroll2019}.

The sequence considered is
\[
  \boxed{\text{Endpoint}\longrightarrow\text{Triangle}
  \longrightarrow\text{Loop}\longrightarrow\text{Curvature}.}
\]
The endpoint map defines the quotient. With composition, a basic
comparison is the direct route $p\to r$ versus the composite comparison
$p\to q\to r$.
When a return route is available, completing that comparison produces
transport around a loop. In smooth
transport geometry, infinitesimal transport around an infinitesimal loop
is represented by curvature.

Quotient topology, category theory, gauge theory, differential geometry,
and general relativity have different structures. The claim here is not
that these theories are derived from endpoint data. It is only that, in
each setting, information depending on a representative, morphism, path,
transport, or loop is determined by endpoint data only when it is
constant on the relevant endpoint fibers.

\section{Endpoint Fibers and Factorization}
\label{sec:bottleneck}

A map from finer descriptions to coarser data may be a quotient, an
identification, a projection to boundary data, or the passage from routes
to endpoint data. The quantities that descend along the map are exactly
the quantities constant on its fibers.

Let $X$ be a space of descriptions and let $q:X\to Q$ be such a map.
A quantity $F:X\to S$ descends to $Q$ precisely when there exists
$\bar F:Q\to S$ such that
\[
  F=\bar F\circ q.
\]
Equivalently, $F$ is constant on the fibers of $q$. This is the ordinary
quotient descent criterion \citep{Munkres2000}. The proof below is a
structured use of that fact.

For comparison, the relevant map is the endpoint map. Routes may
contain more information than endpoints. Endpoint comparison consists only
of the endpoint class. The relevant obstruction is the endpoint fiber: the
class of all routes with identical endpoint data.

\begin{definition}[Endpoint comparison data]
\label{def:endpoint-data}
Endpoint comparison data consist of a class of endpoints $\calP$ and,
for each ordered pair $(p,q)$ for which comparison is allowed, a datum
$C(p,q)$. A route from $p$ to $q$ is an additional object denoted
$\gamma:p\to q$. Two routes have the same endpoint data when they have
the same source, the same target, and the same endpoint comparison datum.
\end{definition}

\begin{definition}[Comparison framework with endpoint data]
\label{def:comparison-framework}
A comparison framework with endpoint data is a tuple
\[
  \mathfrak C=(\calP,\calR,s,t,A,C,\Data),
\]
where $\calP$ is a class of endpoints, $\calR$ is a class of routes,
$s,t:\calR\to\calP$ assign source and target, $A\subseteq\calP\times\calP$
is the class of admissible endpoint pairs, $\Data$ is a class of
endpoint-comparison values, and $C:A\to\Data$ assigns endpoint data. We
assume $(s(\gamma),t(\gamma))\in A$ for every $\gamma\in\calR$. The
endpoint map is
\[
  e:\calR\longrightarrow A\times\Data,
  \qquad
  e(\gamma)=\bigl((s(\gamma),t(\gamma)),C(s(\gamma),t(\gamma))\bigr).
\]
The endpoint equivalence relation is
\[
  \gamma\sim_E\eta \quad\Longleftrightarrow\quad e(\gamma)=e(\eta),
\]
and $\pi_E:\calR\to\calR/{\sim_E}$ is the quotient map. We write
$[\gamma]_E$ for the $\sim_E$-equivalence class of $\gamma$.
\end{definition}

\begin{definition}[Endpoint recovery and constancy on fibers]
\label{def:recoverability}
Let $I:\calR\to S$ be a route invariant. It is determined by endpoint
data, or recoverable from endpoint data, if there exists a function
$J:e(\calR)\to S$ such that $I=J\circ e$.
It is constant on endpoint fibers if, whenever $\gamma\sim_E\eta$, one
has $I(\gamma)=I(\eta)$. It is route-dependent if it distinguishes at least
two routes with the same source, target, and endpoint datum. A framework
includes enough information for $I$ when it contains the distinctions
needed to evaluate $I$.
\end{definition}

\begin{definition}[Endpoint closure for an invariant]
\label{def:endpoint-closure}
Let $I:\calR\to S$ be a route invariant. The endpoint quotient is
\emph{closed for $I$} if $I$ descends to the endpoint quotient; that is,
if there exists a necessarily unique function
$\bar I:\calR/{\sim_E}\to S$ such that
\[
  I=\bar I\circ\pi_E.
\]
A failure of this condition is an \emph{endpoint-closure obstruction} for
$I$.
\end{definition}

\begin{proposition}[Factorization criterion for endpoint data]
\label{prop:factorization}
For a route invariant $I:\calR\to S$, the following are equivalent:
\begin{enumerate}
\item $I$ is determined by endpoint data;
\item the endpoint quotient is closed for $I$;
\item $I$ is constant on endpoint fibers;
\item there exists a necessarily unique function
$\bar I:\calR/{\sim_E}\to S$ such that
\[
  I=\bar I\circ\pi_E.
\]
\end{enumerate}
\end{proposition}

\begin{proof}
Items (2) and (4) are identical by \cref{def:endpoint-closure}. If
$I=J\circ e$, then $\gamma\sim_E\eta$ implies
$e(\gamma)=e(\eta)$ and hence $I(\gamma)=I(\eta)$, so $I$ is constant on
endpoint fibers. If $I$ is
constant on endpoint fibers, define $\bar I([\gamma]_E)=I(\gamma)$. This
is well-defined precisely because all representatives of $[\gamma]_E$
have the same $I$-value, and then $I=\bar I\circ\pi_E$. Finally, if
$I=\bar I\circ\pi_E$, define $J(e(\gamma))=\bar I([\gamma]_E)$. This is
well-defined because equal endpoint data give the same endpoint class,
and then $I=J\circ e$.
\end{proof}

\begin{theorem}[Endpoint-closure criterion]
\label{thm:endpoint-ceiling}
For a route invariant $I:\calR\to S$, the following are equivalent:
\begin{enumerate}
\item $I$ is recoverable from endpoint data;
\item the endpoint quotient is closed for $I$;
\item $I$ factors through the endpoint quotient $\pi_E$;
\item $I$ is constant on endpoint fibers.
\end{enumerate}
Consequently, any route invariant that distinguishes routes with the
same source, target, and endpoint datum is an endpoint-closure
obstruction.
\end{theorem}

\begin{proof}
The equivalences are exactly \cref{prop:factorization} together with
\cref{def:endpoint-closure}. If a route invariant distinguishes
routes in the same endpoint fiber, then it is not constant on that fiber.
It therefore cannot factor through endpoint data, so endpoint closure for
that invariant fails.
\end{proof}

\begin{remark}[Local meaning of endpoint closure]
\label{rem:recoverability-closure}
The term endpoint closure is used here only in the quotient-descent
sense. Given a projection $q:X\to Q$ and a quantity $F:X\to S$, say
that $Q$ is closed for $F$ when $F$ is recoverable from $Q$, equivalently
when $F=\bar F\circ q$ for some $\bar F:Q\to S$. If $q$ is
surjective, such $\bar F$ is unique. Under the specialization
$X=\calR$, $Q=\calR/{\sim_E}$, $q=\pi_E$, and $F=I$,
\cref{thm:endpoint-ceiling} is exactly this quotient-recovery criterion.
The map $e:\calR\to e(\calR)$ has the same fibers as $\pi_E$, so
recovery from endpoint data and descent to the endpoint quotient are
equivalent. Thus endpoint closure is a relative property of one map and
one invariant. It is not a claim about physically closed systems,
dynamics, or completeness of a state space.
All claims in this paper follow from quotient descent for the endpoint
projection.
\end{remark}

\begin{proposition}[Endpoint closure as quotient recovery]
\label{prop:closure-hierarchy}
Let $q:X\to Q$ be a quotient or projection and let $F:X\to S$ be a
quantity. Recovery of $F$ along $q$ is the condition that there exists
$\bar F:Q\to S$ such that $F=\bar F\circ q$. Endpoint closure is
exactly the special case
\[
  (X,Q,q,F)=(\calR,\calR/{\sim_E},\pi_E,I).
\]
Thus an endpoint-closure obstruction is exactly failure of the
factorization $I=\bar I\circ\pi_E$. When the invariant is transport and
reversed routes are available, this obstruction is equivalently detected
by a loop with nontrivial transport; in smooth transport geometries, the
infinitesimal form of local loop transport is curvature. No completeness
or additional closure condition is assumed.
\end{proposition}

\begin{proof}
The first statement is the definition of recovery by factorization.
Substituting $X=\calR$, $Q=\calR/{\sim_E}$, $q=\pi_E$, and $F=I$
yields \cref{def:endpoint-closure}. Failure of endpoint closure is
therefore failure of this factorization. The loop test is
\cref{thm:route-loop}; the curvature statement is the smooth transport
interpretation in \cref{subsec:geometry-gr}.
\end{proof}

\begin{principle}[Criterion for projection to endpoint data]
\label{prin:comparison-compression}
Endpoint comparison data determine exactly the invariants constant on
endpoint fibers. Equivalently, endpoint projection is closed exactly for
those invariants. Every invariant that varies inside an endpoint fiber
requires additional route or transport structure.
\end{principle}

The endpoint quotient is the relevant factorization target.
Before quotienting there may be paths, histories, morphisms, transports,
representatives, or loops. After quotienting there are endpoint classes.
Only information constant on endpoint fibers factors through. In this
sense, only such information is closed under endpoint projection.

\section{Composition, Loops, and Obstruction}
\label{sec:defects}

Composition provides a finite test of whether endpoint data suffice.
With a direct route $p\to r$ and a composite route $p\to q\to r$, the
endpoint pair is the same but the route data may differ. The triangle is
a simple compositional case in which endpoint data may have discarded
route information.

\begin{definition}[Transport comparison structure]
\label{def:transport-structure}
A transport comparison structure consists of a class of endpoints
$\calP$, a class of routes $\calR$ with source and target maps
$s,t:\calR\to\calP$, identity routes $1_p:p\to p$, a partially defined
associative composition $\eta\circ\gamma:p\to r$ whenever
$\gamma:p\to q$ and $\eta:q\to r$, data spaces $E_p$ over endpoints, and
maps
\[
  \tau_\gamma:E_p\to E_q
\]
for every route $\gamma:p\to q$, such that
\[
  \tau_{1_p}=\id_{E_p},\qquad
  \tau_{\eta\circ\gamma}=\tau_\eta\circ\tau_\gamma.
\]
When reversed routes are admitted, each route $\gamma:p\to q$ has a
route $\gamma^{-1}:q\to p$ with
$\tau_{\gamma^{-1}}=\tau_\gamma^{-1}$.
\end{definition}

\begin{proposition}[Endpoint data do not determine route comparison]
\label{prop:endpoint-not-route}
Pairwise endpoint data alone do not determine whether direct comparison
from $p$ to $r$ agrees with comparison through $q$.
\end{proposition}

\begin{proof}
It suffices to exhibit two transport structures with the same endpoint
data and different direct-versus-routed values. Let the endpoint data be
fixed and constant on all admissible pairs. Take endpoints $p,q,r$ and
routes $a:p\to q$, $b:q\to r$, a direct route $c:p\to r$, and the
composite route $d=b\circ a:p\to r$. In both structures let every fiber
be the same two-element set $E=\{0,1\}$ and set
$\tau_a=\tau_b=\id_E$, hence $\tau_d=\tau_b\circ\tau_a=\id_E$. In
the first structure set $\tau_c=\id_E$, so $\tau_c=\tau_d$. In the
second let $\tau_c$ be the transposition of $E$, so $\tau_c\neq\tau_d$.
The endpoint assignment is the same in both structures, but the
direct-versus-composite comparison has a different truth value. Hence
that truth value is not determined by endpoint data alone.
\end{proof}

\begin{corollary}[Triangle criterion]
\label{cor:triangle-obstruction}
If a framework distinguishes the direct route $p\to r$ from the
composite comparison $p\to q\to r$, then the distinguished quantity is not
determined by endpoint data alone.
\end{corollary}

\begin{proof}
The direct route and the composite route have the same source and
target, hence the same endpoint datum $C(p,r)$ under the endpoint map.
If an invariant distinguishes them, it is nonconstant on an endpoint
fiber. The conclusion follows from \cref{thm:endpoint-ceiling}.
\end{proof}

\begin{principle}[Direct-versus-composite comparison]
\label{prin:triangular-first-test}
The triangle is a simple compositional test of projection to comparison
data: direct route versus composite route. A nontrivial distinction
between the two is already a failure of descent to endpoint data.
\end{principle}

\begin{proposition}[Recovery of transport from endpoint data is path independence]
\label{prop:transport-route-invariant}
In a transport comparison structure, view
$\gamma\mapsto\tau_\gamma$ as a map to the disjoint union of all
transport-map sets. It is determined by endpoint data if and only if
transport is path independent: for all routes $\gamma,\eta:p\to q$ with
the same source and target, one has $\tau_\gamma=\tau_\eta$.
\end{proposition}

\begin{proof}
Endpoint data group routes with the same source and target. The invariant
$\gamma\mapsto\tau_\gamma$ factors through endpoint data exactly when it
is constant on each family of routes with the same source and target.
\end{proof}

When a return route is available, the triangle can be completed to a loop.
A loop has degenerate endpoint data: source and target are both $p$. If
transport around the loop is not the identity, the resulting invariant is
not determined by the endpoint pair $(p,p)$.

\begin{definition}[Transport around a loop and holonomy]
\label{def:loop-defect}
For a loop $\gamma:p\to p$ in a transport comparison structure, define
the transport around the loop by
\[
  \delta_p(\gamma)=\tau_\gamma\in\Aut(E_p),
\]
when $\tau_\gamma$ is invertible. The holonomy at $p$ is
\[
  \Hol_p=\{\tau_\gamma\in\Aut(E_p):\gamma:p\to p\text{ is a loop}\}.
\]
The transport around the loop is nontrivial if
$\tau_\gamma\neq\id_{E_p}$.
\end{definition}

\begin{theorem}[Loop test for endpoint-closure obstruction]
\label{thm:route-loop}
Assume transports compose and invert along reversed routes. Let
\[
  S_\tau=\bigsqcup_{p,q}\operatorname{Map}(E_p,E_q),
\]
and let $I:\calR\to S_\tau$ be given by $I(\gamma)=\tau_\gamma$.
Then the following are equivalent:
\begin{enumerate}
\item endpoint closure fails for $I$;
\item $I$ is not recoverable from endpoint data;
\item transport is route-dependent;
\item there exists a loop with nontrivial transport.
\end{enumerate}
Thus, when reversed routes are available, every endpoint-closure
obstruction for transport has a loop test.
\end{theorem}

\begin{proof}
The equivalence of (1), (2), and (3) follows from
\cref{thm:endpoint-ceiling,prop:transport-route-invariant}. Here
endpoint data are the source--target data of the transport structure. It
remains to compare (3) and (4). If transport is route-dependent, choose
routes
$\gamma_1,\gamma_2:p\to q$ with the same source and target and with
$\tau_{\gamma_1}\neq\tau_{\gamma_2}$. The loop
$\gamma=\gamma_2^{-1}\circ\gamma_1:p\to p$ has transport
\[
  \tau_\gamma=\tau_{\gamma_2}^{-1}\circ\tau_{\gamma_1}.
\]
If this were $\id_{E_p}$, then $\tau_{\gamma_1}=\tau_{\gamma_2}$, a
contradiction. Conversely, a loop $\gamma:p\to p$ with
$\tau_\gamma\neq\id_{E_p}$ differs from the constant loop $1_p$ while
sharing the same endpoint data. Thus transport is route-dependent.
\end{proof}

\begin{remark}[Scope of the loop test]
The equivalence with a nontrivial loop is asserted in the reversible
transport regime stated in Theorem~\ref{thm:route-loop}. Without
reversed routes, the endpoint-fiber criterion remains the controlling
statement: endpoint closure fails exactly when the invariant is
nonconstant on an endpoint fiber.
\end{remark}

\begin{corollary}[Holonomy is not determined by endpoint data]
\label{cor:holonomy-endpoint}
Any framework with nontrivial holonomy contains holonomy information not
recoverable from endpoint data alone.
\end{corollary}

\begin{proof}
A nontrivial loop $\gamma:p\to p$ and the constant loop $1_p$ have the
same endpoint data but different transports. Hence holonomy is
nonconstant on an endpoint fiber and is not determined by endpoint data.
\end{proof}

Thus the chain is a sequence of tests. Projection to endpoint
data defines the quotient. Composition tests endpoint closure. The triangle
is a simple compositional test. A loop is the case with identical source
and target.
Holonomy is loop information not determined by endpoint data.

\section{Standard Structures Used}
\label{sec:existing-frameworks}

The factorization result is the quotient factorization statement from
\cref{prop:factorization}. The cited references provide standard
interpretations of routes, endpoint quotients, and route invariants.

The quotient fact is standard: a function descends through a quotient
exactly when it is constant on equivalence classes \citep{Munkres2000}.
In category theory the corresponding data are objects, morphisms,
identities, and associative composition
\citep{EilenbergMacLane1945,MacLane1998}. In differential geometry the
corresponding data are connections, parallel transport, holonomy, and
curvature; Wilson loops are closed-loop gauge observables; the
Ambrose--Singer theorem relates holonomy to curvature
\citep{Ehresmann1950,Wilson1974,AmbroseSinger1953,KobayashiNomizu1963,Nakahara2003}.
In general relativity the corresponding hierarchy is metric--connection--
curvature--contraction \citep{Einstein1916,Wald1984,Malament2012,Carroll2019}.

In each case, the relevant test is to specify the routes, specify the
endpoint quotient, and check constancy of the invariant on endpoint
fibers.

\begin{proposition}[External instance schema]
\label{prop:external-schema}
Suppose a cited theory contains a class of routes, an endpoint-data
quotient on those routes, and a route invariant $I$ distinguishing some
routes with the same endpoint data. Then the failure of $I$ to be
determined by endpoint data follows formally from
\cref{thm:endpoint-ceiling}.
\end{proposition}

\begin{proof}
Interpret $\calR$, $\sim_E$, and $I$ in the cited theory. If there are
routes $\gamma\sim_E\eta$ with $I(\gamma)\neq I(\eta)$, then $I$ is
nonconstant on a fiber of $\pi_E$. By \cref{prop:factorization}, $I$
cannot factor through endpoint data.
\end{proof}

The cited ingredients used here are quotient descent
\citep{Munkres2000}; category
language with distinct parallel arrows
\citep{EilenbergMacLane1945,MacLane1998}; connections \citep{Ehresmann1950};
closed-loop gauge observables \citep{Wilson1974}; the curvature--holonomy
relation and differential-geometric background
\citep{AmbroseSinger1953,KobayashiNomizu1963}; holonomy/path structures in
general relativity and Yang--Mills theory \citep{Barrett1991}; and the
relativistic ordering from metric to connection to curvature and the Einstein
tensor \citep{Wald1984,Malament2012,Carroll2019}.

\section{Instances Where Endpoint Data Do Not Determine the Invariant}
\label{sec:realizations}

In each case there is a domain-specific invariant. The criterion above
states the condition under which that invariant fails to factor through
endpoint data.

\subsection{Quotients and Categorical Morphisms}
\label{subsec:quotients-categories}

A quotient groups representatives. A quantity descends to the
quotient exactly when it is independent of the chosen representative
\citep{Munkres2000}.

\begin{proposition}[Quotient descent as recovery from endpoint data]
\label{prop:quotient-descent}
Let $q:X\to X/{\sim}$ be a quotient map. A function $F:X\to S$ is
determined by quotient data if and only if there exists
$\bar F:X/{\sim}\to S$ such that $F=\bar F\circ q$. Hence any
representative-dependent $F$ is not determined by quotient data.
\end{proposition}

\begin{proof}
The statement is the quotient descent criterion. If $F=\bar F\circ q$,
then equivalent representatives have the same value. Conversely, if $F$
is constant on equivalence classes, define $\bar F([x])=F(x)$. This is
well-defined exactly because $F$ is constant on each equivalence class.
\end{proof}

The categorical version is compositional. Objects are endpoints;
morphisms are routes; parallel morphisms are routes with the same source
and target. Object data do not determine morphism-dependent information
\citep{EilenbergMacLane1945,MacLane1998}.

\begin{proposition}[Object data do not determine morphism invariants]
\label{prop:category-realization}
Let $\calC$ be a category and let
$I:\Mor(\calC)\to S$ be a morphism invariant. Let
$o:\Mor(\calC)\to\Ob(\calC)^2$ be the map sending
a morphism to its ordered source--target pair. Then $I$ is determined
by object data alone, meaning $I=J\circ o$ for some $J$, if and only if
$I$ is constant on every hom-class $\Hom_{\calC}(X,Y)$.
In particular, any invariant distinguishing parallel morphisms
$f,g:X\to Y$ is not determined by object data.
\end{proposition}

\begin{proof}
For fixed objects $X,Y$, the corresponding fiber of $o$ is
$\Hom_{\calC}(X,Y)$. The result is the same fiber-constancy criterion as
\cref{prop:factorization}.
\end{proof}

\subsection{Gauge Transport and Network Path Products}
\label{subsec:gauge-networks}

In gauge theory, two identifications are easily confused. Gauge
invariance removes dependence on representative. It does not remove
dependence on route. Wilson loops are gauge-invariant loop data, not
endpoint data \citep{Wilson1974}. Connections supply the transport along
paths \citep{Ehresmann1950,Nakahara2003}.

\begin{proposition}[Nontrivial gauge holonomy is not determined by endpoint data]
\label{prop:gauge-realization}
Let a gauge theory assign parallel transport $U_\gamma$ to paths
$\gamma:p\to q$, and let endpoint data group paths with the same source
and target. If a gauge-invariant quantity $I$ built from the
transport, for example a Wilson-loop value after completing paths to loops,
has different values on two paths with the same source and target, then
$I$ is not determined by endpoint data alone.
\end{proposition}

\begin{proof}
The two paths lie in one endpoint fiber. Since $I$ has different values
on them, it is not constant on endpoint fibers. By
\cref{prop:factorization}, it cannot factor through endpoint data.
\end{proof}

Networks provide a finite model of the criterion. Vertices are
endpoints. Edges carry labels. Paths multiply labels. No independent
network theorem is used here; the example is a finite instance of the
factorization criterion above.

\begin{proposition}[Network path products are determined by endpoint data exactly
under path independence]
\label{prop:network-realization}
Let a directed network have edge labels in a monoid $M$, and let
$P(\gamma)$ be the ordered product of labels along a finite path
$\gamma$, with $P(1_p)$ equal to the unit of $M$. Let endpoint data on
paths be their ordered source--target pair. Then $P$ is determined by
endpoint data if and only if all paths with the same source and target
have the same product. In particular, any nontrivial loop product
$P(\gamma)\neq P(1_p)$ is not determined by endpoint data.
\end{proposition}

\begin{proof}
The endpoint fibers are exactly the families of paths with the same
source and target. By
\cref{prop:factorization}, the path product factors through endpoints
exactly when it is constant on those fibers.
\end{proof}

\subsection{Differential Geometry and General Relativity}
\label{subsec:geometry-gr}

A connection geometry provides a smooth instance of the same factorization
problem. Points are endpoints. Tangent spaces or bundle fibers are data
spaces. A connection is the rule for comparing fibers along paths
\citep{Ehresmann1950,KobayashiNomizu1963,Nakahara2003}. Holonomy is
transport around loops. According to the Ambrose--Singer theorem, the
infinitesimal generators of holonomy are curvature data transported back
to the base point \citep{AmbroseSinger1953,KobayashiNomizu1963}.
Barrett treats holonomy/path structures in general relativity and
Yang--Mills theory \citep{Barrett1991}.

In general relativity the same order appears in Lorentzian form. The
metric provides pointwise comparison. Compatible torsion-free transport is
Levi-Civita transport. Curvature is formed from that connection. Ricci
and scalar contractions enter the Einstein tensor
\citep{Einstein1916,Wald1984,Malament2012,Carroll2019}. The hierarchy is
not endpoint--curvature, but endpoint comparison, transport, loop
obstruction, curvature, and contraction.

\begin{proposition}[Nontrivial holonomy is not determined by endpoint data]
\label{prop:curvature-realization}
In a smooth connection geometry, if parallel transport has nontrivial
holonomy, then the corresponding holonomy invariant is not determined by
endpoint data. For local or restricted holonomy, the standard
Ambrose--Singer curvature--holonomy relation determines the holonomy algebra
from curvature data transported back to the base point. Curvature is
therefore the smooth infinitesimal form of this failure of descent to
endpoint data.
This statement does not identify every global holonomy effect with
curvature.
\end{proposition}

\begin{proof}
A nontrivial loop $\gamma:p\to p$ and the constant loop $1_p$ have the
same endpoint data but different transports. Thus holonomy is not
determined by endpoint data. For local or restricted holonomy, the
Ambrose--Singer theorem determines the holonomy algebra from curvature data
transported back to the base point.
Thus curvature is the infinitesimal version of this route-dependent
level; global flat holonomy is a separate route-dependent obstruction.
\end{proof}

Curvature records the infinitesimal transport obstruction; global flat
holonomy records a non-infinitesimal route obstruction.

\section{General Endpoint-Closure Obstruction}
\label{sec:thesis}

In each example, the obstruction is nonconstancy on endpoint fibers.

\begin{theorem}[Endpoint-closure obstruction]
\label{thm:general-extension}
Let a comparison framework contain endpoint data, composable route data,
and a route invariant $I$ capable of distinguishing routes with the same
endpoint data. Then endpoint closure fails for $I$. Equivalently, $I$ is
not determined by endpoint data, $I$ is not constant on endpoint fibers,
and $I$ does not factor through the endpoint quotient $\pi_E$.
Consequently, endpoint comparison data determine only invariants constant
on endpoint fibers; representation of $I$ requires additional route or
transport structure.
\end{theorem}

\begin{proof}
If $I$ distinguishes routes $\gamma$ and $\eta$ with the same endpoint
data, then $\pi_E(\gamma)=\pi_E(\eta)$ but $I(\gamma)\neq I(\eta)$. Thus
$I$ is nonconstant on an endpoint fiber. By
\cref{thm:endpoint-ceiling}, this is exactly failure of endpoint closure
for $I$, equivalently failure of factorization through $\pi_E$.
\end{proof}

\begin{proposition}[Preservation under fiber-preserving translation]
\label{prop:structural-universality}
Let $\mathfrak C$ and $\mathfrak D$ be comparison frameworks with route
classes $\calR_{\mathfrak C}$ and $\calR_{\mathfrak D}$ and endpoint
equivalence relations $\sim_E^{\mathfrak C}$ and $\sim_E^{\mathfrak D}$.
Let $T:\calR_{\mathfrak C}\to\calR_{\mathfrak D}$ be a translation such
that
\[
  \gamma\sim_E^{\mathfrak C}\eta
  \quad\Longrightarrow\quad
  T(\gamma)\sim_E^{\mathfrak D}T(\eta).
\]
If route invariants $I_{\mathfrak C}$ and $I_{\mathfrak D}$ satisfy
$I_{\mathfrak C}=I_{\mathfrak D}\circ T$, and if endpoint closure fails
for $I_{\mathfrak C}$, then endpoint closure fails for
$I_{\mathfrak D}$ on the image of $T$.
\end{proposition}

\begin{proof}
Since endpoint closure fails for $I_{\mathfrak C}$, by
\cref{thm:endpoint-ceiling} there are routes
$\gamma\sim_E^{\mathfrak C}\eta$ with
$I_{\mathfrak C}(\gamma)\neq I_{\mathfrak C}(\eta)$. The fiber condition
on $T$ implies $T(\gamma)\sim_E^{\mathfrak D}T(\eta)$, while preservation
of invariant values implies
$I_{\mathfrak D}(T(\gamma))\neq I_{\mathfrak D}(T(\eta))$. Hence
$I_{\mathfrak D}$ is nonconstant on an endpoint fiber of the translated
image. By \cref{thm:endpoint-ceiling}, endpoint closure fails there.
\end{proof}

Gauge holonomy, categorical morphism dependence, network path products,
and Riemann curvature are distinct structures. In each case, the same
formal endpoint-closure obstruction is present: nonconstancy on endpoint
fibers.

\section{Scope of the Criterion}
\label{sec:scope}

The endpoint-closure criterion is not a derivation of connections,
gauge fields, metrics, curvature tensors, or equations of motion. It is
a limitation result: endpoint data are too coarse whenever the relevant
invariant is not constant on endpoint fibers. Representation of that
invariant requires additional domain-specific structure.

Nor are dynamics derived. Gauge equations, variational
principles, conservation laws, Einstein's field equations, and empirical
interpretation require assumptions not present in endpoint data.
In general relativity, curvature is downstream of transport. It does not
imply the Einstein equations.

The relation among these examples is formal rather than explanatory:
each case has an invariant that fails to descend along the relevant
endpoint projection. The shared structure is nonconstancy on endpoint
fibers, not an additional closure axiom.

Nontrivial holonomy is not a matter of endpoint convention under this
criterion. If a smooth geometry that includes the invariant has
nontrivial holonomy, endpoint data do not suffice for the relevant
invariant. For local or restricted holonomy, curvature is the
infinitesimal form of that nonconstancy; global flat holonomy remains a
separate route-dependent obstruction.

\section{Conclusion}
\label{sec:conclusion}

The conclusion is the following:
\[
\boxed{\text{Endpoint closure holds exactly for invariants constant on fibers.}}
\]
Equivalently, invariants determined by endpoint data are exactly the
invariants that factor through the quotient by endpoint-data equivalence.
Information separating routes with the same source and target is
nonconstant on endpoint fibers. Its representation requires additional
route or transport structure.

The sequence is:
\[
  \boxed{\text{Endpoint}\longrightarrow\text{Triangle}
  \longrightarrow\text{Loop}\longrightarrow\text{Curvature}.}
\]
The endpoint map defines the quotient. A direct route together with a
composite route forms the triangle. A return route completes the triangle
to a loop. In smooth transport
geometry, infinitesimal transport around an infinitesimal loop is
represented by curvature.

Failure of quotient descent, categorical morphism dependence, network path
dependence, gauge holonomy, Wilson loops, differential-geometric
holonomy, curvature, and the curvature hierarchy of general relativity
all concern the same formal question: which invariants fail to
factor through endpoint data, and hence obstruct endpoint closure?

The obstruction is route-dependent information. Endpoint descriptions
are limited by the endpoint map. For any framework using only endpoint
data, the condition is constancy on endpoint fibers, equivalently endpoint
closure for the invariant under consideration.

\begin{sloppypar}
\raggedright
\begin{thebibliography}{99}

\bibitem{AmbroseSinger1953}
Ambrose, W., \& Singer, I.~M. (1953).
A theorem on holonomy.
\textit{Transactions of the American Mathematical Society}, \textit{75}(3),
428--443. https://doi.org/10.1090/S0002-9947-1953-0063739-1

\bibitem{Barrett1991}
Barrett, J.~W. (1991).
Holonomy and path structures in general relativity and Yang--Mills theory.
\textit{International Journal of Theoretical Physics}, \textit{30}(9),
1171--1215. https://doi.org/10.1007/BF00671007

\bibitem{Carroll2019}
Carroll, S.~M. (2019).
\textit{Spacetime and geometry: An introduction to general relativity}
(2nd ed.). Cambridge University Press. https://doi.org/10.1017/9781108770385

\bibitem{Ehresmann1950}
Ehresmann, C. (1951).
Les connexions infinit\'esimales dans un espace fibr\'e diff\'erentiable.
In \textit{Colloque de Topologie (Bruxelles, 1950)} (pp.~29--55).
Centre Belge de Recherches Math\'ematiques.

\bibitem{EilenbergMacLane1945}
Eilenberg, S., \& Mac~Lane, S. (1945).
General theory of natural equivalences.
\textit{Transactions of the American Mathematical Society}, \textit{58},
231--294. https://doi.org/10.1090/S0002-9947-1945-0013131-6

\bibitem{Einstein1916}
Einstein, A. (1916).
Die Grundlage der allgemeinen Relativit\"atstheorie.
\textit{Annalen der Physik}, \textit{354}(7), 769--822.
https://doi.org/10.1002/andp.19163540702

\bibitem{KobayashiNomizu1963}
Kobayashi, S., \& Nomizu, K. (1963).
\textit{Foundations of differential geometry} (Vol.~I).
Wiley-Interscience.

\bibitem{MacLane1998}
Mac~Lane, S. (1998).
\textit{Categories for the working mathematician} (2nd ed.).
Springer. https://doi.org/10.1007/978-1-4757-4721-8

\bibitem{Malament2012}
Malament, D.~B. (2012).
\textit{Topics in the foundations of general relativity and Newtonian
gravitation theory}. University of Chicago Press.

\bibitem{Munkres2000}
Munkres, J.~R. (2000).
\textit{Topology} (2nd ed.). Prentice Hall.

\bibitem{Nakahara2003}
Nakahara, M. (2003).
\textit{Geometry, topology and physics} (2nd ed.).
Institute of Physics Publishing.

\bibitem{Wald1984}
Wald, R.~M. (1984).
\textit{General relativity}. University of Chicago Press.

\bibitem{Wilson1974}
Wilson, K.~G. (1974).
Confinement of quarks.
\textit{Physical Review D}, \textit{10}(8), 2445--2459.
https://doi.org/10.1103/PhysRevD.10.2445

\bibitem{WolfeCSM}
Wolfe, C.~K. (2026).
\textit{Closed Systems from Comparison Completeness: A Theory of Closed
Relational Systems}.
Closure Research Initiative monograph, Version 7.
\url{https://closureresearchinitiative.org/csm/}

\end{thebibliography}
\end{sloppypar}

\end{document}
